\documentclass[11pt]{article}

\usepackage[margin=1in]{geometry}
\usepackage{amsmath,amssymb,bm}
\usepackage{microtype}
\usepackage[numbers,sort&compress]{natbib}
\usepackage[colorlinks=true,allcolors=blue]{hyperref}
\usepackage[nameinlink,noabbrev]{cleveref}

\newcommand{\mbf}[1]{\mathbf{#1}}
\newcommand{\tr}{\operatorname{tr}}
\newcommand{\dd}{\,\mathrm d}
\newcommand{\pder}[2]{\frac{\partial #1}{\partial #2}}
\providecommand{\doi}[1]{\href{https://dx.doi.org/#1}{#1}}

\title{A constrained variational reduction for a finite-deformation\
single-solid, single-fluid Biot system}
\author{Exploratory manuscript}
\date{}

\begin{document}
\maketitle

\begin{abstract}
The mineral equation in a finite-deformation Biot model relates mineral-volume
ratio, pore pressure, and elastic volume.  We show that this equation is the
local stationarity condition of a pressure-dependent mineral potential.  This
observation gives a constrained variational interpretation of the
finite-deformation Biot stress transformation, but it does not by itself
remove a global equation: mineral volume is already a locally condensed
constitutive state.  For a nonreacting solid with compatible referential mass,
solid mass conservation provides a distinct continuum reduction.  It eliminates
the global solid partial-density field and gives a two-field system for
displacement and pore pressure, with mineral volume still found by a scalar
local solve.  We derive the reduced storage, total stress, weak equations,
and the conditions under which the reduction is valid.  Darcy flow requires
a dissipative virtual-power or Onsager extension; it does not follow from a
conservative Hamilton principle alone.  The result identifies a focused
two-field prototype while retaining the constitutive mineral response that
produces the nonlinear Biot coefficient.
\end{abstract}

\noindent\textbf{Keywords:} finite-deformation poromechanics; constrained
variation; Biot coefficient; solid mass conservation; Darcy dissipation

\section{Introduction}
\label{sec:introduction}

In a saturated porous material, total volume change and mineral-volume change
are different kinematic responses.  Their distinction determines the solid
volume fraction, fluid storage, and the contribution of pore pressure to the
mixture stress.  A finite-deformation Biot coefficient follows from the
mineral-volume response at fixed pore pressure \citep{fosterxu2025}.  The
matched logarithmic specialization considered here supplies that response by
a scalar equation coupling elastic volume, pressure, and mineral volume.

The scalar equation appears naturally as a constraint candidate.  This paper
asks two separate questions.  First, can that relation be obtained from a
variational statement?  Second, does making it explicit reduce the number of
global equations?  The answers differ.  The relation is a local stationarity
condition of a mixed mineral potential.  Its elimination is therefore an
energy-consistent constitutive condensation.  It does not reduce the existing
global system, because the mineral-volume ratio is already local rather than
a finite-element field.  A separate reduction follows from source-free solid
mass conservation: for a solid with prescribed compatible referential mass, it eliminates
the solid partial density in favor of the deformation.

The contribution is a precise separation of these two eliminations for the
single-solid model, together with the resulting continuum two-field weak
problem.  It is a specialization and reformulation of the existing model,
rather than a claim that a mineral multiplier supplies a new field reduction.

We restrict attention to one deformable solid and one compressible fluid at
constant temperature.  There are no solid sources, reactions, or phase
changes.  The construction follows the constituent-volume distinction used
in mixture theories \citep{drumheller2000} and retains the finite-deformation
pressure transformation of \citet{fosterxu2025}.  Section~\ref{sec:state}
states the model.  Section~\ref{sec:stationarity} derives the local
stationarity result and the total stress.  Section~\ref{sec:reduction}
performs the global reduction.  Section~\ref{sec:variational-flow} separates
the conservative and dissipative portions of the formulation.

\section{State, mass balance, and mineral equation}
\label{sec:state}

Let \(\mbf x=\boldsymbol\chi(\mbf X,t)\) be the solid motion,
\(\mbf F=\nabla_{\mbf X}\boldsymbol\chi\), and \(J=\det\mbf F>0\).
Let \(\bar J\) denote the mineral-volume ratio and let \(\phi_{s0}\) be
the reference solid volume fraction.  The solid partial density \(\rho_s\)
is mass per current mixture volume, whereas the solid intrinsic density
\(\bar\rho_s\) is mass per current mineral volume.  Suppose the material
solid has no mass source and its initial data satisfy
\(J\rho_s|_{t=0}=\rho_{s0}(\mbf X)\).  Integrating solid mass balance then
gives
\begin{align}
 J\rho_s &= \rho_{s0}(\mbf X),
 \label{eq:solid-mass-integrated}\\
 \bar J &= \frac{\bar\rho_{s0}}{\bar\rho_s},
 \qquad
 \phi_s=\frac{\phi_{s0}\bar J}{J}.
 \label{eq:phase-volume-state}
\end{align}
The constant-reference specialization uses
\(\rho_{s0}=\phi_{s0}\bar\rho_{s0}\).  Thus \(\bar J\) determines the
intrinsic density, while the ratio \(\phi_{s0}\bar J/J\) determines the
current solid volume fraction.  A spatially varying prescribed
\(\rho_{s0}(\mbf X)\) is also admissible, but the formulas below use the
constant-reference specialization.

For the elastic specialization, let \(J^e\) be elastic volume.  It equals
\(J\) in the absence of plastic distention and equals \(J/a^p\) when a
stored plastic volume factor \(a^p\) is held fixed.  The matched logarithmic
mineral relation is
\begin{equation}
 R(J^e,p,\bar J)
 =K_s\ln\bar J
 +\left(1-\frac{K}{\phi_{s0}K_s}\right)p\bar J
 -\frac{K}{\phi_{s0}}\ln J^e=0,
 \label{eq:mineral-residual}
\end{equation}
where \(K\) and \(K_s\) are the drained and mineral reference bulk moduli,
with \(0<K<\phi_{s0}K_s\).  Write
\(\alpha=1-K/(\phi_{s0}K_s)>0\).  We follow the root continued from the
drained state and retain it where
\begin{equation}
 D\equiv\pder{R}{\bar J}
 =\frac{K_s}{\bar J}
 +\left(1-\frac{K}{\phi_{s0}K_s}\right)p>0.
 \label{eq:stable-mineral-branch}
\end{equation}
Equivalently, \(K_s+\alpha p\bar J>0\).  This condition makes the chosen
scalar branch regular; its strict-local-minimum meaning follows below.

The finite-deformation Biot coefficient is defined at fixed pressure,
referential solid mass, and plastic state by
\begin{equation}
 B=1-\phi_{s0}\left.\pder{\bar J}{J}\right|_{p,J\rho_s,\mbf F^p}.
 \label{eq:biot-definition}
\end{equation}
Implicit differentiation of \cref{eq:mineral-residual} gives
\begin{equation}
 B=1-
 \frac{K\bar J}
 {J\left[K_s+\left(1-\dfrac{K}{\phi_{s0}K_s}\right)p\bar J\right]}.
 \label{eq:biot-coefficient}
\end{equation}
At \(J=\bar J=1\) and \(p=0\), this recovers \(B=1-K/K_s\)
\citep{biotwillis1957}.

\section{Mineral equation as a local stationarity condition}
\label{sec:stationarity}

For the elastic specialization, the relevant part of the matched free energy
per unit reference mixture volume is
\begin{equation}
 W_s(\mbf F^e,\bar J)
 =W_{\mathrm{iso}}(\mbf F^e)
 +\frac K2(\ln J^e)^2
 +\frac{\phi_{s0}K_s}{2\alpha}
 \left(\ln\bar J-\frac{K}{\phi_{s0}K_s}\ln J^e\right)^2,
 \label{eq:matched-energy}
\end{equation}
where \(W_{\mathrm{iso}}\) is independent of \(\bar J\).  In the
following constitutive variations, the plastic state is held fixed, so that
\(\mbf F^e\) is determined by \(\mbf F\).  The pressure conjugacy is
\begin{equation}
 p=-\frac{1}{\phi_{s0}}
 \left.\pder{W_s}{\bar J}\right|_{\mbf F^e}.
 \label{eq:pressure-conjugacy}
\end{equation}
Substitution of \cref{eq:matched-energy} gives
\begin{equation}
 \pder{W_s}{\bar J}+\phi_{s0}p
 =\frac{\phi_{s0}}{\alpha\bar J}R(J^e,p,\bar J).
 \label{eq:energy-residual-identity}
\end{equation}
For prescribed \((\mbf F,p)\), introduce
\begin{equation}
 \Pi_s(\mbf F,\bar J;p)=W_s(\mbf F,\bar J)+\phi_{s0}p\bar J.
 \label{eq:local-mineral-potential}
\end{equation}
Equations \cref{eq:energy-residual-identity,eq:mineral-residual} show that
\(\delta_{\bar J}\Pi_s=0\) is exactly the mineral equation.  At a root,
its second variation is
\begin{equation}
 \frac{\partial^2\Pi_s}{\partial\bar J^2}
 =\frac{\phi_{s0}}{\alpha\bar J}
 \pder{R}{\bar J}
 =\frac{\phi_{s0}\left(K_s+\alpha p\bar J\right)}
 {\alpha\bar J^2}>0.
 \label{eq:mineral-curvature}
\end{equation}
Hence the selected positive-derivative root is a strict local minimum and a
regular implicit branch.  This statement neither establishes a global minimum
nor excludes other roots under tensile pressure.

The same result supplies the stress transformation.  Define the mixed
mechanical potential per unit reference volume
\begin{equation}
 G(\mbf F,\bar J,p)
 =W_s(\mbf F,\bar J)-p\left(J-\phi_{s0}\bar J\right).
 \label{eq:mixed-mechanical-potential}
\end{equation}
Let \(W^{\prime\prime}(\mbf F,p)=W_s(\mbf F,\bar J(\mbf F,p))\) and
\(\mbf P^{\prime\prime}=\partial W^{\prime\prime}/\partial\mbf F|_p\).
Writing \(\mbf P'=\partial W_s/\partial\mbf F|_{\bar J}\), the fixed-
pressure chain rule and \cref{eq:pressure-conjugacy} give
\begin{align}
 \mbf P^{\prime\prime}
 &=\mbf P'-\phi_{s0}p
 \left.\pder{\bar J}{J}\right|_pJ\mbf F^{-T},
 \label{eq:fixed-pressure-stress-chain-rule}\\
 &=\mbf P'-(1-B)pJ\mbf F^{-T}.
 \label{eq:prime-stress-transform}
\end{align}
The envelope theorem applied to \cref{eq:mixed-mechanical-potential} now
gives the total derivative along the stationary state,
\begin{equation}
 \frac{\dd}{\dd\mbf F}G(\mbf F,\bar J(\mbf F,p),p)
 =\mbf P'-pJ\mbf F^{-T}
 =\mbf P^{\prime\prime}-BpJ\mbf F^{-T}.
 \label{eq:mixed-stress-total}
\end{equation}
This is a fixed-pressure mechanical identity.  The potential \(G\) is not,
by itself, an action for a free pressure variation: its derivative with
respect to \(p\) contains pore volume rather than the fluid mass balance.

An explicit multiplier constraint \(\lambda R=0\) is therefore unnecessary
for this model.  It introduces an additional local variable \(\lambda\) and
an additional equation.  Eliminating them recovers the same local tangent
obtained by differentiating the scalar mineral solve.  The stationarity form
explains the existing condensation; it does not create a further global
reduction.

\section{Continuum reduction of the source-free system}
\label{sec:reduction}

The continuum reduction of a global equation comes instead from
\cref{eq:solid-mass-integrated}.  The original three-field representation
uses \((\mbf u,p,\rho_s)\), with \(\mbf u=\mbf x-\mbf X\).  Because
\(\rho_s=\rho_{s0}/J\) follows pointwise from the integrated solid mass
balance and the compatible initial condition, the solid partial-density field
and its residual can be removed at the continuum level.
The locally solved mineral ratio remains necessary:
\begin{equation}
 R(J^e,p,\bar J)=0,
 \qquad
 \phi_s=\frac{\phi_{s0}\bar J}{J}.
 \label{eq:reduced-local-state}
\end{equation}

Let \(\bar\rho_f(p)\) be a barotropic fluid intrinsic density and let
\(\mbf W_f\) be fluid mass flux relative to the solid, expressed per unit
reference area.  The referential fluid mass becomes
\begin{equation}
 m_f(\mbf F,p)
 =J(1-\phi_s)\bar\rho_f(p)
 =\left(J-\phi_{s0}\bar J(\mbf F,p)\right)\bar\rho_f(p).
 \label{eq:reduced-fluid-storage}
\end{equation}
Let \(\Omega_0\) be the solid reference domain, with displacement and
traction boundaries \(\Gamma_u\) and \(\Gamma_t\), and pressure and normal
mass-flux boundaries \(\Gamma_p\) and \(\Gamma_w\).  Let \(\mbf b_0\) be
body force per reference volume.  The two-field governing equations are
\begin{align}
 \nabla_{\mbf X}\mathbin{\cdot}
 \left(\mbf P^{\prime\prime}-BpJ\mbf F^{-T}\right)
 +\mbf b_0&=\mbf 0,
 \label{eq:reduced-momentum}\\
 \dot m_f+\nabla_{\mbf X}\mathbin{\cdot}\mbf W_f&=0.
 \label{eq:reduced-water-mass}
\end{align}
For isotropic permeability \(\kappa\), viscosity \(\mu_f\), and zero
gravity, Darcy flow is
\begin{equation}
 \mbf W_f=-\bar\rho_f J\mbf F^{-1}
 \frac{\kappa}{\mu_f}\mbf F^{-T}\nabla_{\mbf X}p.
 \label{eq:darcy-flux}
\end{equation}
The state dependency in \cref{eq:reduced-fluid-storage} includes the
implicit response of \(\bar J\) to both deformation and pressure.  Thus the
reduced water equation retains the mineral-compressibility contribution to
storage and the same Biot coefficient in momentum.

For test functions \(\mbf w\) and \(q\), vanishing on \(\Gamma_u\) and
\(\Gamma_p\), respectively, the corresponding weak residuals are
\begin{align}
 \mathcal R_u(\mbf w)
 &=\int_{\Omega_0}\nabla_{\mbf X}\mbf w:
 \left(\mbf P^{\prime\prime}-BpJ\mbf F^{-T}\right)\dd V_0
 -\int_{\Omega_0}\mbf w\mathbin{\cdot}\mbf b_0\dd V_0
 -\int_{\Gamma_t}\mbf w\mathbin{\cdot}\bar{\mbf t}\dd A_0,
 \label{eq:reduced-weak-momentum}\\
 \mathcal R_p(q)
 &=\int_{\Omega_0}q\dot m_f\dd V_0
 -\int_{\Omega_0}\nabla_{\mbf X}q\mathbin{\cdot}\mbf W_f\dd V_0
 +\int_{\Gamma_w}q\bar w_f\dd A_0.
 \label{eq:reduced-weak-water}
\end{align}
This is a two-field continuum system only under the stated source-free,
compatible-initial-data assumptions.  Solid production, reactions, or evolving
referential mass restore a solid mass equation and its field.  The reduction
also removes freedom to prescribe \(\rho_s\) independently of the initial
motion and \(\rho_{s0}(\mbf X)\).

The result does not assert algebraic equivalence to a finite-element method
that independently approximates \(\rho_s\).  For example, substituting
\(\rho_{s0}/J(\mbf u_h)\) into a Q2--Q1--Q2 discretization does not generally
produce a Q2 density field.  A prototype should use displacement and pressure
spaces selected for the reduced storage, evaluate \(\bar J\) at quadrature
points, and differentiate the complete local state.  It should then compare
pressure and displacement with the three-field solution, check referential
solid and fluid mass drift under time refinement, and compare its assembled
Jacobian with finite differences.  The coupled PDE in this paper is restricted
to elasticity.  A poroplastic extension requires a local return map and the
plastic-state contribution to \(\dot m_f\) and the outer tangent.

\section{Conservative variation and Darcy dissipation}
\label{sec:variational-flow}

The stationary mineral calculation and the mechanical virtual work have a
conservative variational structure.  For quasistatic mechanics, arbitrary
admissible \(\delta\mbf u\) gives the internal virtual work
\begin{equation}
 \delta\mathcal W_{\mathrm{int}}
 =\int_{\Omega_0}
 \left(\mbf P^{\prime\prime}-BpJ\mbf F^{-T}\right):
 \nabla_{\mbf X}\delta\mbf u\dd V_0.
 \label{eq:mechanical-virtual-work}
\end{equation}
Together with external virtual work, this yields
\cref{eq:reduced-momentum}.  With inertia, the same internal variation
enters an action in the usual Hamilton form.

Fluid transport is different: Darcy flux dissipates energy.  A conservative
Hamilton action alone cannot return \cref{eq:darcy-flux}.  It must be
supplemented by a dissipative virtual-power term or an Onsager rate
minimization \citep{bedford1979variational,doi2011onsager}.  Let
\(\mu_f^{\mathrm{chem}}\) denote the fluid chemical potential per unit mass.
At constant temperature,
\(\nabla_{\mbf X}\mu_f^{\mathrm{chem}}=\nabla_{\mbf X}p/\bar\rho_f\).
The
mass-flux rate potential and force pairing are
\begin{equation}
 \mathcal O(\mbf W_f)
 =\frac{\mu_f}{2\bar\rho_f^2\kappa J}
 \mbf W_f\mathbin{\cdot}
 \left(\mbf F^T\mbf F\right)\mbf W_f.
 \label{eq:darcy-onsager-potential}
\end{equation}
Stationarity of the Onsager functional
\begin{equation}
 \mathcal O(\mbf W_f)+\mbf W_f\mathbin{\cdot}
 \nabla_{\mbf X}\mu_f^{\mathrm{chem}}
 \label{eq:darcy-onsager-functional}
\end{equation}
with respect to \(\mbf W_f\) gives
\begin{equation}
 \frac{\mu_f}{\bar\rho_f^2\kappa J}
 \left(\mbf F^T\mbf F\right)\mbf W_f
 +\frac{1}{\bar\rho_f}\nabla_{\mbf X}p=\mbf0,
 \label{eq:darcy-onsager-stationarity}
\end{equation}
which is \cref{eq:darcy-flux}.  Fluid mass balance remains the kinematic
constraint that relates this flux to the storage in
\cref{eq:reduced-fluid-storage}.  A complete time-incremental mixed
principle must include that mass constraint and a fluid free energy; the
fixed-pressure mechanical potential \(G\) alone is not such a principle.

This separation identifies the appropriate variational program for the
two-field prototype.  The mineral equation should be condensed locally from
\cref{eq:local-mineral-potential}; the reduced fluid storage should use
\cref{eq:reduced-fluid-storage}; and Darcy flow should enter through a
Hamilton--d'Alembert or Onsager dissipation term.  Treating the mineral
equation as a global multiplier constraint would obscure, rather than
simplify, this structure.

\section{Conclusions}
\label{sec:conclusions}

The finite-deformation mineral equation is an energy stationarity condition
in the mineral-volume ratio.  Its selected strict-local-minimum branch permits a local,
energy-consistent elimination of mineral volume and produces the usual
finite-deformation total stress.  That elimination is already present when a
material routine solves the scalar mineral equation, so an explicit
constraint multiplier does not reduce the global system.

For the nonreacting single-solid, single-fluid specialization, source-free
solid mass conservation yields a different continuum reduction.  It removes the
solid partial-density field, leaving displacement and pressure as global
unknowns while preserving the local mineral solve, nonlinear fluid storage,
and Biot stress response.  A future implementation should compare this
two-field system against the existing three-field formulation under matching
initial and boundary data, then determine whether the reduction retains the
desired robustness in drainage and finite-deformation benchmarks.

\bibliographystyle{plainnat}
\bibliography{explorations/constrained_variational_biot/references}
\end{document}
